\subsection{Kinematics}
The throw trajectory of the robot arm was calculated using cubic polynomials with specified joint angles and velocities at the start and end points. The calculated trajectories follow a linear movement in joint space, but a non-linear movement in cartesian space. Alternatively, the trajectory could follow a linear movement in cartesian space but a non-linear movement in joint space. Planning a linear movement in joint space makes it easier to find a trajectory to satisfy the joint space safety limits (table~\ref{tab:jointLimits}). However, planning a linear movement in cartesian space makes it easier to find a trajectory to satisfy the safety planes in cartesian space (table~\ref{tab:tcpLimits}). As we were given safety limitations in both the joint space and cartesian space, this made the trajectory planning particularly challenging. In the end we chose to use a linear movement in joint space. This was because the joint limits were the most restrictive compared to the cartesian limits, especially joint 5 which has only 40 degrees acceptable range. On top of that, trying to make a linear movement in cartesian space using the “cubicpolytraj” \cite{cubicPolyTraj} function resulted in the tcp making a curved movement between the specified points instead of connecting them in a straigth line. Therefore, it would be necessary to manually write a function if had chosen to make a linear movement in cartesian space. Another method to deal with the joint restrictions would be to lock the most problematic joint and do the kinematics without it. This would also require removing this joint from the Jacobian matrix when calculating the joint velocities \cite{kinematicsLecture10}. However, this approach was not chosen in the end, as not using all joints made it more difficult to find inverse kinematics solution for various transformation matrices. Finally, we chose cubic polynomials over quintic polynomials as it was not necessary to specify acceleration constraints, only velocity and position constraints. Alternatively, the cubic polynomials could be combined into a single trapezoidal velocity profile with the required joint velocities for the release point in the middle and 0 rad/s in the start and end for every joint. This is something that could be experimented in the future to further develop this project \cite{kinematicsLecture9}.

\subsection{Evaluation}
The tests for the final evaluation could be improved and expanded. First, instead of a square box, it would be better to use a round box. This is because for the square box the distance between the center and edge is larger in the corners than between the corners, while a round box has a constant radius all around. Additionally, instead of only testing 11 throws with random target coordinates leading to a single success rate, multiple targets could be tested systematically. This could be done by testing multiple points spread systematically over the table space with a fixed distance between them. The target box would be placed at each target point and multiple throws would be tested for each point. This would yield a success rate for each tested point across the table space. Using this data a 2D heat map of success rates across the table space could be created. From this data, a total average success rate over the whole tested space could be calculated. The same approach could also be taken for the accuracy test to create a heat map of the throw accuracy over the table space and to calculate the average accuracy over the whole tested space. 
